%------------------------------------------------------------------------------%
\section{Demonstração por contraposição}

Sabemos que as proposições
\[
  P \Rightarrow Q \quad\text{ e }\quad \sim Q \Rightarrow \sim P
\]
são equivalentes, ou seja,
\[
  P \Rightarrow Q \quad\Leftrightarrow\quad \sim Q \Rightarrow \sim P.
\]
Nesse caso então iremos demonstrar de forma direta a sentença
$\sim Q \Rightarrow \sim P$,
que implica na demonstração de $P \Rightarrow Q$. Vamos
demonstrar as proposições que seguem:

%------------------------------------------------------------------------------%
\begin{proposicao}{Número é par se seu quadrado for par}{prop:par_quadrado}
  Se $x^2$ é par, então $x$ é par.
\end{proposicao}
%------------------------------------------------------------------------------%

%------------------------------------------------------------------------------%
\begin{proof}
  Inicialmente é importante observar, que nesse caso fica difícil fazer a
  demonstração direta (Verifique!). Vamos demonstrar a veracidade dessa
  afirmação usando a contrapositiva. Nesse caso, temos:
  \[
    P\colon x^2 \text{ é par}.
  \]
  \[
    Q\colon x \text{ é par}.
  \]
  Sendo assim,
  \[
    \sim Q\colon x \text{ é ímpar}.
  \]
  \[
    \sim P\colon x^2 \text{ é ímpar}.
  \]
  Assumindo que $x$ é ímpar, temos que $x=2k+1$, $k \in \Z$. Assim,
  \[
    x^2
    = (2k+1)^2
    = 4k^2+4k+1
    = 2(2k^2+2k)+1
  \]
  Chamando $p=2k^2+2k$ e notando que $2k^2+2k$ é um número inteiro, temos que
  $x^2=2p+1$, que é o mesmo que dizer que $x^2$ é ímpar.
\end{proof}

